\begin{center}
  \CachedTikzFigure{content/resources/chapters/real_functions/tikz/functions_symmetry.tex}
\end{center}

\par\noindent\begin{minipage}{\linewidth}
\begin{DefinitionBox}{Funktionseigenschaft -- Gerade und ungerade Funktionen}{real_functions_parity}
  Sei $D$ symmetrisch zu null, also $x\in D\Rightarrow -x\in D$.
  Eine Funktion $f:D\to\R$ heißt
  \begin{itemize}
    \item \emph{gerade} \EnglishTerm{even}, wenn
      $f(-x)=f(x)$ für alle $x\in D$ gilt.
      Ihr Graph ist achsensymmetrisch zur $y$-Achse.
    \item \emph{ungerade} \EnglishTerm{odd}, wenn
      $f(-x)=-f(x)$ für alle $x\in D$ gilt.
      Ihr Graph ist punktsymmetrisch zum Ursprung.
  \end{itemize}
  \par\smallskip
  \begin{minipage}[t]{0.49\linewidth}
    \centering
    \textcolor{THMBlue}{Gerade Funktion}
    \par\smallskip
    \CachedTikzFigure{content/resources/chapters/real_functions/tikz/functions_even_symmetry_schematic.tex}
  \end{minipage}\hfill
  \begin{minipage}[t]{0.49\linewidth}
    \centering
    \textcolor{THMBlue}{Ungerade Funktion}
    \par\smallskip
    \CachedTikzFigure{content/resources/chapters/real_functions/tikz/functions_odd_symmetry_schematic.tex}
  \end{minipage}
\end{DefinitionBox}
\begin{ExampleBox}[before skip=0.3\baselineskip]{Symmetrie rechnerisch prüfen}{real_functions_parity_check}
  Wir setzen jeweils $-x$ ein und vergleichen den entstehenden Term mit
  $f(x)$ und $-f(x)$.
  Untersuchen Sie dazu die drei Funktionen auf $\R$:
  \[
    f(x)=x^4-2x^2,\qquad g(x)=x^3-x,\qquad h(x)=x^2+x.
  \]
  \begin{ClozeArea}[34]{9}
    \[
      \begin{array}{l@{\quad}l@{\quad}l}
        f(x)=x^4-2x^2&
        f(-x)=x^4-2x^2=f(x)&\text{gerade},\\[0.8em]
        g(x)=x^3-x&
        g(-x)=-x^3+x=-g(x)&\text{ungerade},\\[0.8em]
        h(x)=x^2+x&
        h(-x)=x^2-x&\text{weder noch}.
      \end{array}
    \]
    Bei $h$ genügt bereits $x=1$ als Gegenbeispiel:
    $h(-1)=0$ ist weder $h(1)=2$ noch $-h(1)=-2$.
  \end{ClozeArea}
\end{ExampleBox}
\end{minipage}


\begin{KeyPointBox}{Symmetrie spart Arbeit}{real_functions_parity_use}
  Bei geraden und ungeraden Funktionen genügt es, eine Hälfte des Graphen
  zu untersuchen. Die andere Hälfte folgt durch Spiegelung beziehungsweise
  Drehung um $180^\circ$.

  Ist eine ungerade Funktion bei $x=0$ definiert, gilt notwendig
  $f(0)=0$. Die Umkehrung gilt nicht: Eine Nullstelle bei null allein
  beweist keine Symmetrie.
\end{KeyPointBox}

\begin{TheoremBox}{Symmetrie beim Rechnen mit Funktionen}{real_functions_parity_rules}
  Die folgende Übersicht zeigt, wie sich gerade und ungerade Funktionen
  beim Rechnen verhalten:
  \begin{center}
    \renewcommand{\arraystretch}{1.55}
    \begin{tabularx}{\linewidth}{|>{\raggedright\arraybackslash}p{2.5cm}||>{\raggedright\arraybackslash}X|>{\raggedright\arraybackslash}X|}
      \hline
      & \textcolor{THMBlue}{$g$ gerade}
      & \textcolor{green!55!black}{$g$ ungerade}\\ \hline\hline
      \textcolor{THMBlue}{$f$ gerade}
      & $f(x)\pm g(x)$, $f(x)g(x)$, $f(x)/g(x)$ und $f(g(x))$ sind
        \textcolor{THMBlue}{gerade}.
      & $f(x)g(x)$ und $f(x)/g(x)$ sind \textcolor{green!55!black}{ungerade};
        $f(g(x))$ ist \textcolor{THMBlue}{gerade}.\\ \hline
      \textcolor{green!55!black}{$f$ ungerade}
      & $f(x)g(x)$ und $f(x)/g(x)$ sind \textcolor{green!55!black}{ungerade};
        $f(g(x))$ ist \textcolor{THMBlue}{gerade}.
      & $f(x)\pm g(x)$ und $f(g(x))$ sind
        \textcolor{green!55!black}{ungerade};
        $f(x)g(x)$ und $f(x)/g(x)$ sind \textcolor{THMBlue}{gerade}.\\ \hline
    \end{tabularx}
  \end{center}
  Die Aussagen gelten auf geeigneten, zu null symmetrischen
  Definitionsmengen; für $f/g$ muss zusätzlich $g(x)\ne0$ gelten.
  Gemischte Summen sind im Allgemeinen weder gerade noch ungerade.
\end{TheoremBox}
Die Produktregel für zwei ungerade Funktionen folgt direkt durch Einsetzen:
\[
  (fg)(-x)=f(-x)g(-x)=(-f(x))(-g(x))=f(x)g(x)=(fg)(x).
\]
